% ECONOMICS_PROBLEM: {"id": "O34-212", "title": "Innovation rents and competition", "jel_code": "O34", "source_id": "OP-0268"}
% !TeX program = xelatex
\documentclass[11pt,a4paper]{article}
\usepackage[margin=23mm]{geometry}
\usepackage{fontspec}
\setmainfont{Latin Modern Roman}
\usepackage{amsmath,amssymb,mathtools}
\usepackage{xcolor,enumitem,booktabs,longtable,array,xurl,hyperref}
\definecolor{muted}{HTML}{53636C}
\hypersetup{colorlinks=true,urlcolor=blue,linkcolor=blue}
\setlength{\parindent}{0pt}
\setlength{\parskip}{5pt}
\newcommand{\R}{\mathbb R}
\newcommand{\E}{\mathbb E}
\newcommand{\N}{\mathbb N}
\newcommand{\ind}{\mathbf 1}
\newcommand{\Var}{\operatorname{Var}}
\newcommand{\Cov}{\operatorname{Cov}}
\newcommand{\supp}{\operatorname{supp}}
\DeclareMathOperator*{\argmin}{arg\,min}
\DeclareMathOperator*{\argmax}{arg\,max}
\newcommand{\entrymeta}[3]{{\sffamily\small\color{muted}\textbf{\texttt{#1}}\quad|\quad #2\par #3\par}}
\newcommand{\Problemref}[1]{\texttt{#1}}
\newcommand{\JELref}[2]{\texttt{#2} (JEL \texttt{#1})}
\begin{document}
\section*{Innovation rents and competition}
\textbf{Problem ID:} \texttt{O34-212}\quad \textbf{JEL:} \texttt{O34}\par
% BEGIN ECONOMICS STATEMENT
\label{problem:OP-0268}
{\sffamily\small\color{muted}Original subject group: Growth and productivity\par}
\label{definitions:problem:30007549}
\textbf{Audit classification:} Explicit published open question. Competition and innovation rents are expressly a policy challenge; a compact patent/entry optimum is editorial. Verification concerns the cited version, not first priority or proof that this exact editorial specification remains unsolved on October 8, 2026.
\subsection*{Definitions and precise research target}
\paragraph{Editorial mathematical specification.} Fix a Schumpeterian product-line economy in which incumbents and entrants choose research effort and successful innovations improve quality by a specified amount. Patent protection grants innovation rents for duration \(T\), while competition policy \(a\) determines a declared entry restriction or merger rule. Researchers are scarce, consumers have specified preferences, and diffusion produces an explicitly modeled knowledge spillover. Let \(\mathcal P\) be a compact set of feasible pairs \((T,a)\), excluding unspecified discretionary interventions.
For an initial state and discount factor \(\beta\in(0,1)\), define social welfare
\[
W(T,a)=E_{T,a}\sum_{t\ge0}\beta^t
u(C_t).
\]
Here \(C_t\) is aggregate consumption after research and enforcement resource use, and \(u\) is household utility; these costs enter resource feasibility rather than being subtracted twice. Characterize the maximizing feasible policy and identify the response parameters needed to rank alternatives. Include induced entry, incumbent research, replacement, diffusion, and equilibrium research wages; treating innovation rates as fixed while changing rents misses the central tradeoff. Determine which protections reward socially valuable innovation and which mainly preserve obsolete incumbents. Validate predicted innovation responses using credible policy variation and distinguish rent measurement from technological quality. Report how the optimum changes with spillovers, enforcement, and privately known quality. This is a conditional dynamic policy-design target, not a universal assertion that stronger competition always raises or lowers innovation.

A Schumpeterian economy here means that successful innovations improve a product's quality and may replace its incumbent producer; specify line states, quality-step distribution, entry costs, research-intensity technologies, labor supply, and final-demand aggregation as primitives. Let \(T\) be patent duration in periods and \(a\) a policy parameter with a defined mapping into feasible entry/merger decisions. Under each \((T,a)\), firms optimize discounted profit, research labor clears, consumers optimize, and the common selected equilibrium determines \(C_t\). Require \(\sum_t\beta^tE|u(C_t)|<\infty\); explosive growth with utility growing too fast can invalidate the objective. Compact policies do not establish attainment without continuity and nonempty equilibria. The target is a policy ranking conditional on the entire primitive family and welfare weights, not a theorem giving a universal competition optimum.
\subsection*{Variants and assumption changes}
The following are \emph{editorial variants}, not additional published open conjectures.
\begin{enumerate}
\item \emph{Patent duration only.} Hold entry/merger policy fixed and vary patent expiration with a stated postexpiration licensing rule. Compare incentives to invent with subsequent diffusion and enforcement resource use.
\item \emph{Acquisition policy.} Hold intellectual-property rules fixed and model incumbent acquisition of potential entrants. Compare a merger restriction with unrestricted acquisitions, including changes in entrants' expected exit value and their research effort.
\end{enumerate}
\subsection*{References and source audit}
\begin{itemize}
\item Ufuk Akcigit. \emph{Creative destruction in economic growth}. 2026. Locator: Section7, Open questions, item3. Primary text checked. \url{https://onlinelibrary.wiley.com/doi/10.1111/sjoe.70037}.
\end{itemize}
Competition and innovation rents are expressly a policy challenge; a compact patent/entry optimum is editorial.
\begin{enumerate}\raggedright
\item\hypertarget{definitions:source:30007549:1}{} Ufuk Akcigit. 2026. \emph{Creative destruction in economic growth}. Section7, Open questions, item3.
\url{https://onlinelibrary.wiley.com/doi/10.1111/sjoe.70037}.
DOI: \url{https://doi.org/10.1111/sjoe.70037}.
\end{enumerate}

\subsection*{Common conventions: Source mathematical conventions}

These conventions apply to each entry unless explicitly replaced.

\textbf{Models and identification.} A primitive model \(m\) specifies all probability laws, feasible choices, information, equilibrium and measurement rules used by its target. The admissible class \(\mathcal M\) is fixed independently of outcomes. If \(P_O(m)\) is its observable probability law and \(\tau(m)\) the target, its identified set at observable law \(P\) is
\[
 \mathcal I_\tau(P)=\{\tau(m):m\in\mathcal M,\ P_O(m)=P\}.
\]
Point identification means that this set is a singleton. An empty set signals incompatibility of data and assumptions. Coordinate bounds need not describe the joint set. Bounds are sharp only if every retained target is attainable or arbitrarily approachable by compatible models. Finite bounds require explicit restrictions on unobserved quantities; unrestricted confounding, hidden wealth, extrapolation or arbitrary equilibrium selection can yield uninformative sets.

\textbf{Causal interventions.} A potential outcome \(Y_i(p)\) is the outcome under a complete policy regime \(p\), including timing, financing, information and market spillovers. The target population and units are fixed before treatment. With interference, use the complete assignment vector or a justified exposure mapping; individual no-interference notation alone cannot identify an economy-wide intervention. A difference of mean potential outcomes does not require the cross-world joint distribution. Conditioning on post-treatment migration, survival, enrollment or employment changes the estimand and needs separate justification.

\textbf{Statistical statements.} An estimator, confidence set, forecast or test specifies a sampling law, model class and loss/coverage criterion. Uniform \((1-\alpha)\)-coverage means \(\inf_{m\in\mathcal M}\Pr_m\{\tau(m)\in C_n\}\geq1-\alpha\), or an explicitly asymptotic version. The whole parameter space trivially has coverage, so informativeness requires a stated width, risk or power objective. A forecasting improvement is relative to a named benchmark, information set and evaluation population.

\textbf{Equilibrium and optimization.} An equilibrium comprises feasible plans, optimality, beliefs where needed, budget constraints and all market/resource clearing. The primitives or a stated benchmark define each correspondence \(\mathcal E(p;m)\). Nonemptiness is assumed or proved, never inferred just from compact policy sets. A selected-equilibrium target uses a declared selection rule; a robust target quantifies over all permitted equilibria. Use suprema or infima unless attainment is established. An \(\varepsilon\)-optimal policy is feasible and within \(\varepsilon\) of the finite optimum value. Report an empty feasible set separately.

\textbf{Welfare and accounting.} Normative utility, interpersonal normalization, social weights, numeraire, discounting and the use of public receipts are primitives. Behavioral utility need not coincide with the welfare criterion. Household welfare includes ownership income and rents once; adding profits again to compensating variation generally double-counts them. A monetary equivalent is defined by an explicit transfer and price convention, with monotonicity and range conditions. Average effects, mechanism contributions, transfers and welfare losses are different targets. Infinite sums require convergence and dynamic feasible plans require terminal or no-Ponzi conditions.

\textbf{Source versions.} A working-paper date, revision date and journal publication year are distinguished. A DOI for a journal version is not automatically the DOI of an earlier manuscript. Locators refer to the stated version and printed pages where specified. A metadata-only check does not verify a claim about a full-text conclusion. Every entry's source subsection narrows the provenance claim accordingly.
% END ECONOMICS STATEMENT
\end{document}
